\chapter{Referencias y material de partida}
\label{sec:referencias}

\begin{enumerate}[label={[\arabic*]}]
  \item Muñoz Muñoz, Alfonso. \textit{Criptografía aplicada y Esteganografía. Tarea evaluable}. Máster en Ciberseguridad. Documento facilitado para la práctica: \texttt{tarea.pdf}, p.~3.
  \item Material docente del módulo: documentación y diapositivas de los temas 1.1 a 1.9.
  \item OpenSSL Project. Sitio web indicado en el enunciado. \url{https://www.openssl.org/}.
  \item GeeksforGeeks. \textit{Practical Uses of OpenSSL Command in Linux}. Referencia indicada en el enunciado. \url{https://www.geeksforgeeks.org/practical-uses-of-openssl-command-in-linux/}.
  \item\label{ref:tema-kdf} Muñoz Muñoz, Alfonso. \textit{Tema 1.4. Funciones criptográficas de resumen HASH y XOF}, versión 1.0.0. Material docente del módulo, pp.~9--19. XOF (Extendable-Output Function). Derivación, sales, PBKDF2 (pp.~12--14) y composición Encrypt-then-MAC.
  \item\label{ref:slides-kdf} Muñoz Muñoz, Alfonso. \textit{Módulo 1.4. Funciones criptográficas de resumen y XOF}. Diapositivas del módulo, diap.~7, 21 y 26--27. HMAC, Encrypt-then-MAC y derivación de claves.
  \item\label{ref:tema-aes} Muñoz Muñoz, Alfonso. \textit{Tema 1.2. Algoritmos criptográficos simétricos. Cifrado en bloque y flujo}, versión 1.0.0. Material docente del módulo, pp.~6, 15, 21 y 24. AES, CBC, relleno y uso de OpenSSL.
  \item\label{ref:rfc-pbkdf2} Moriarty, K., Kaliski, B. y Rusch, A. \textit{PKCS \#5: Password-Based Cryptography Specification Version 2.1}. RFC~8018 (Request for Comments), enero de 2017, secciones 4.1 y 5.2. \url{https://www.rfc-editor.org/rfc/rfc8018.html}.
  \item\label{ref:openssl-pbkdf2}\label{ref:openssl-kdf-cli} OpenSSL Project. \textit{EVP\_KDF-PBKDF2} y \textit{openssl-kdf}. Documentación oficial de los parámetros y de la interfaz de línea de comandos. \url{https://docs.openssl.org/3.6/man7/EVP_KDF-PBKDF2/}. \url{https://docs.openssl.org/3.6/man1/openssl-kdf/}.
  \item\label{ref:slides-gcm} Muñoz Muñoz, Alfonso. \textit{Módulo 1.2. Algoritmos criptográficos simétricos. Cifrado en bloque y flujo}. Diapositivas del módulo, diap.~16, 37 y 46. Recomendación de AES-GCM con claves de 128 bits y cifrado autenticado.
  \item\label{ref:tema-gcm} Muñoz Muñoz, Alfonso. \textit{Tema 1.2. Algoritmos criptográficos simétricos. Cifrado en bloque y flujo}, versión 1.0.0. Material docente del módulo, pp.~18 y 22--25. Modo GCM, nonce y ejemplos de cifrado autenticado.
  \item\label{ref:tema-aead} Muñoz Muñoz, Alfonso. \textit{Tema 1.4. Funciones criptográficas de resumen HASH y XOF}, versión 1.0.0. Material docente del módulo, pp.~21--23. AEAD, funcionamiento de GCM y verificación de la etiqueta.
  \item\label{ref:nist-gcm} Dworkin, M. \textit{Recommendation for Block Cipher Modes of Operation: Galois/Counter Mode (GCM) and GMAC}. National Institute of Standards and Technology (NIST), SP~800-38D, secciones 5--8. \url{https://csrc.nist.gov/pubs/sp/800/38/d/final}.
  \item\label{ref:openssl-gcm} OpenSSL Project. \textit{openssl-enc}, sección de modos admitidos, y \textit{EVP\_EncryptInit}, sección de cifrado autenticado. Documentación oficial de las limitaciones de \texttt{enc} y del uso de GCM mediante la biblioteca. \url{https://docs.openssl.org/master/man1/openssl-enc/}. \url{https://docs.openssl.org/master/man3/EVP_EncryptInit/}.
\end{enumerate}

